\section{A native proof of the original regular-variation estimate}
\label{sec:el-variation}
This section supplies a complete proof of
\eqref{eq:v109-regular-variation}, which was stated as inherited in
Input~\ref{input:v109-analytic}. Its two averaged conclusions
\eqref{eq:v109-paired-average}--\eqref{eq:v109-volume-average} are
not proved by this calculation and remain separate inputs. The proof
uses the original guard, including every unvisited obstacle, and the
original insertion $I_m$. Nothing is conditioned on a smaller event.

\begin{lemma}[Logarithmic sensitivity of the positive optical denominator]
\label{lem:el-logden}
Regarding $\mathfrak b_m$ in \eqref{eq:v109-optical-product} as a
polynomial in the positive variables $\tau_i$ and $\kappa_j$, it has
nonnegative coefficients and degree at most one in each variable.
Consequently
\begin{equation}\label{eq:el-logden}
 0\leq\frac{\tau_i\partial_{\tau_i}\mathfrak b_m}{\mathfrak b_m}\leq1,
 \qquad
 0\leq\frac{\kappa_j\partial_{\kappa_j}\mathfrak b_m}{\mathfrak b_m}\leq1.
\end{equation}
\end{lemma}
\begin{proof}
Each variable occurs in exactly one optical factor, that factor is
linear in it, and every entry of every factor has nonnegative
coefficients. Matrix multiplication preserves coefficient positivity
and does not repeat a variable. For any variable $x$, write
$\mathfrak b_m=a+xb$ with $a,b\geq0$ and $\mathfrak b_m>0$.
Then $x\partial_x\mathfrak b_m/\mathfrak b_m=xb/(a+xb)\in[0,1]$.
For $m=1$ this applies to $\mathfrak b_1=\tau_0$ and there are no
curvature variables.
\end{proof}

Set $c_\circ=\sqrt{19}/10$ and define
\begin{equation}\label{eq:el-omega-constant}
 C_\omega=
 \frac{2R_+J_*(1+\sigma)}{\tau_-(1-\sigma)}
       +\frac{4C_c}{c_\circ}+2C_c.
\end{equation}

\begin{proposition}[The precise inherited incidence insertion suffices]
\label{prop:el-logomega}
At every original exact-event regular point and for every endpoint
direction,
\begin{equation}\label{eq:el-logomega}
 |D_q\log\omega[\eta]|
     \leq C_\omega I_m\|\eta\|_\infty,
 \qquad I_m=1+\sum_{i=1}^{m-1}\frac{q_0^{d_i}}{c_i}.
\end{equation}
\end{proposition}
\begin{proof}
The radius is fixed under $D_q$, so
\eqref{eq:v109-omega} and Lemma~\ref{lem:el-logden} imply
\[
 |D\log\omega|
 \leq\frac{|Dc_0|}{c_0}+\frac{|Dc_m|}{c_m}
       +\sum_{i=0}^{m-1}\frac{|D\tau_i|}{\tau_i}
       +\sum_{i=1}^{m-1}\frac{|Dc_i|}{c_i}.
\]
Here $D\log\kappa_i=-D\log c_i$; the sign is irrelevant only after
taking absolute values. The original section cap gives
$c_0,c_m\geq c_\circ$ and $s_0,s_m\leq2$. Also
\[
 \sum_{i=0}^{m-1}(s_i+s_{i+1})
 \leq\frac{2(1+\sigma)}{1-\sigma},\qquad
 s_i\leq2\sigma^{d_i}\leq2q_0^{d_i}.
\]
Apply Proposition~\ref{prop:el-observables}. The first two terms in
\eqref{eq:el-omega-constant} bound the endpoint and flight contributions,
and $2C_c(I_m-1)$ bounds the interior contribution. Since $I_m\geq1$,
\eqref{eq:el-logomega} follows with the displayed constant. No additional
inverse incidence or collision-dependent multiplier occurs.
\end{proof}

\subsection{The actual guard, with nonsmooth clearance switches}
Put $M_\Theta=\sup_{x\in\R}|\Theta'(x)|<\infty$ for the fixed
smooth cutoff already present in \eqref{eq:v76-guard}, and set
\begin{equation}\label{eq:el-guard-constant}
 \beta=\sigma q_0^{-1/4}<1,\qquad
 C_W=\frac{4M_\Theta(C_c+C_\Delta)}{1-\beta}.
\end{equation}

\begin{theorem}[Collision-uniform variation of the original physical guard]
\label{thm:el-guard}
For each $\varepsilon>0$, $W^\varepsilon$ is locally Lipschitz on
an actual regular branch. Along every continuously differentiable
physical branch path it is absolutely continuous on compact parameter
intervals and, almost everywhere,
\begin{equation}\label{eq:el-guard}
 |D_qW^\varepsilon[\eta]|
       \leq C_W\varepsilon^{-1}\|\eta\|_\infty.
\end{equation}
The bound is independent of the collision count, exact labels, and
radius. No section-distance factor is inserted in $W^\varepsilon$.
\end{theorem}
\begin{proof}
All guard factors lie in $[0,1]$. Apply the product rule along a compact
path and use Proposition~\ref{prop:el-observables}; for clearance factors
use the almost-everywhere chain rule for Lipschitz functions on an interval.
The collision and flight depth schedules are exactly those in the
original guard, not new schedules. Each depth occurs at most twice, so
\[
 \sum_{i=0}^{m}\frac{|Dc_i|}{q_0^{d_i/4}}
 \leq\frac{4C_c}{1-\beta}\|\eta\|_\infty,
 \qquad
 \sum_{i=0}^{m-1}\frac{|D\Delta_i|}{q_0^{d'_i/4}}
 \leq\frac{4C_\Delta}{1-\beta}\|\eta\|_\infty.
\]
Multiplication by $M_\Theta/\varepsilon$ proves \eqref{eq:el-guard}.
The inequality $\beta<1$ follows, for example, from $\sigma<q_0<1$.
The clearance is locally Lipschitz even at a nearest-disk switch, so
there is no omitted jump contribution at such a switch. Actual section
boundaries and collision-word boundaries are not crossed by this
argument; their one-sided physical traces retain their earlier status.
\end{proof}

\subsection{Roofwise variation and physical normalization}
Keep $\lambda_\circ=9\sqrt{19}/200$ and $\Lambda_\circ=47/6$.
Define
\begin{equation}\label{eq:el-regular-constant}
 C_{\rm reg}=\frac{\Lambda_\circ}{\lambda_\circ}
       \left(C_\omega+\frac{2C_F}{\lambda_\circ}+C_W\right).
\end{equation}

\begin{theorem}[Native regular variation on every actual roof-angle segment]
\label{thm:el-regular}
On every regular segment of the original physical pairing space,
with $u=F$, $e=\nabla F/|\nabla F|$ fixed along the segment,
$z=|\nabla F|>0$, and
$A=\omega/(4\pi c_*\chi)$, the original density
$f^\varepsilon=AW^\varepsilon$ is locally absolutely continuous and
\begin{equation}\label{eq:el-regular}
 |\partial_u f^\varepsilon|
 \leq C_{\rm reg}(1+\varepsilon^{-1})\frac{AI_m}{z}
 \quad\hbox{almost everywhere.}
\end{equation}
This is precisely \eqref{eq:v109-regular-variation}, with a finite
explicit constant depending only on the original geometry and cutoff.
Its proof uses no pressure, curve-growth, mixing, or local-limit theorem.
\end{theorem}
\begin{proof}
The two-dimensional determinant defining $\chi$ equals
$e^{\perp\mathsf T}Ge^\perp$, as in the retained Gauss-coordinate
calculation. Thus $\chi\geq\lambda_\circ$. Differentiating on a fixed-angle
segment, not in an arbitrary direction with changing $e$, gives
\[
 \left\|\frac{dq}{du}\right\|_\infty
 =\left\|\frac{G^{-1}e}{z e^{\mathsf T}G^{-1}e}\right\|_\infty
 \leq\frac{\Lambda_\circ}{\lambda_\circ z},\qquad
 |\partial_u\chi|\leq2C_F\left\|\frac{dq}{du}\right\|_\infty.
\]
The latter bound is Theorem~\ref{thm:el-third}; there is no angular
derivative term because $e$ is fixed. By Proposition~\ref{prop:el-logomega},
\[
 |\partial_u\log A|
 \leq\left(C_\omega I_m+\frac{2C_F}{\lambda_\circ}\right)
                              \frac{\Lambda_\circ}{\lambda_\circ z}.
\]
The factor $4\pi c_*$ is the unchanged physical probability
normalization and is constant under this differentiation. Using
$0\leq W^\varepsilon\leq1$ and Theorem~\ref{thm:el-guard} now yields
\[
 |\partial_u(AW^\varepsilon)|\leq
 \frac{A\Lambda_\circ}{\lambda_\circ z}
 \left(C_\omega I_m+\frac{2C_F}{\lambda_\circ}
                         +\frac{C_W}{\varepsilon}\right).
\]
Since $I_m\geq1$, this implies \eqref{eq:el-regular} with
\eqref{eq:el-regular-constant}. Compact subsegments have positive
incidences and $z>0$, so all differentiations and absolute-continuity
claims are legitimate there. The uniform constants do not depend on
those compact-subsegment minima.
\end{proof}

\begin{corollary}[Explicit integration up to the actual endpoint]
\label{cor:el-traces}
For two roof values in one original regular fixed-angle segment,
\begin{equation}\label{eq:el-traces}
 |f^\varepsilon(u_2)-f^\varepsilon(u_1)|
 \leq C_H(1+\varepsilon^{-1})\sigma^{(m-2)_+}\sqrt{|u_2-u_1|},
 \quad
 C_H=\frac{500C_{\rm reg}\sqrt{2\Lambda_\circ}}
                  {4\pi c_*\lambda_\circ^2}.
\end{equation}
This supplies a native regular-interior premise for the retained
endpoint-trace results, rather than a new claim of zero trace on a
bulk-null unvisited or section boundary.
\end{corollary}
\begin{proof}
Use the already proved physical trajectory bound
$\omega I_m\leq500\sigma^{(m-2)_+}$ in
Corollary~\ref{cor:v109-marked}, $\chi\geq\lambda_\circ$, and integrate
\eqref{eq:el-regular}. Lemma~\ref{lem:v109-gauss} gives
$\int_{u_1}^{u_2}du/z\leq\lambda_\circ^{-1}
\sqrt{2\Lambda_\circ|u_2-u_1|}$. This proves \eqref{eq:el-traces} on
compact subsegments and then at one-sided limits. Identification of
those limits with the actual physical pairing and treatment of
simultaneous singular endpoints remain the retained trace construction;
no derivative of a grazing collision map is introduced here.
\end{proof}

\begin{proposition}[Complete positive variation domination, before any supremum]
\label{prop:el-positive}
On the retained countable one-copy sheet space, let $h\geq0$ be any
Borel selector (including the original exact record and actual roof).
Then, allowing both sides to be infinite,
\begin{equation}\label{eq:el-positive}
 \int h(u,\theta)|\partial_u f^\varepsilon(u,\theta)|\,du\,d\theta
 \leq C_{\rm reg}(1+\varepsilon^{-1})
       \int h(u,\theta)\frac{A I_m}{z}(u,\theta)\,du\,d\theta.
\end{equation}
For selectors $h$ pulled back from physical states, the right side is
$C_{\rm reg}(1+\varepsilon^{-1})
\int_\Pi h\,I_m z^{-1}\,d\nu_R^*$ under the original coarea convention.
\end{proposition}
\begin{proof}
Multiply \eqref{eq:el-regular} by $h$ on compact regular subsegments,
integrate, then exhaust each open segment and the countable sheet space.
All quantities are nonnegative, so monotone convergence proves
\eqref{eq:el-positive} without a word-count factor or a signed-series
exchange. The final identification is precisely the retained physical
one-copy coarea input. It is not a new construction of pairing
multiplicity and does not identify the paired minimum of endpoint
traces with the interior density.
\end{proof}

\subsection{What is closed, and the unchanged remaining mathematical limit}
The regular-variation clause of Input~\ref{input:v109-analytic} now has
a full native geometric proof. The previous trace, clearance, radial,
and critical coefficient theorems remain in the manuscript and are not
reclassified as new results. In particular the two averaged clauses
of that input are not consequences of bounded $D^3F$. Neither
\eqref{eq:el-positive} nor the per-trajectory exponential factor in
\eqref{eq:el-traces} can be multiplied by a physical-word count to
obtain a normalized roof-density estimate.

The original unfinished assertion is still \eqref{eq:v112-remaining},
namely the same positive low-load shortage after subtraction of the
actual nongrazing capacity, with the complete physical sum and roof
essential supremum preceding the collision limsup, and the outer
limit in $B$ last. A proof must still supply the complete normalized
same-physical-roof source or a valid capacity comparison, with the
actual pairing endpoint measure and the first-order clearance and
incidence-only sources retained. The spectral/arithmetic and
conditional-denominator conclusions retain their exact hypotheses.
This section removes a specific local optical proof dependency; it does
not replace this missing global result by a new residual or a relaxed
window. The full raw pointwise local limit theorem and final
four-coordinate inversion are not asserted as completed.
